% 本文件由 scripts/make_tables.py 自动生成，请勿手工编辑
\begin{table}[htbp]
  \centering
  \caption{4 h 之后环境条件外推假设的灵敏度分析（扰动幅度取自附件 1 末小时的实测上下限）}
  \label{tab:sensitivity}
  \scriptsize
  \begin{tabular}{ccccc}
    \toprule
    扰动（相对末值保持） & 问题 2、3 终点/h & 问题 2、3 变化/h & 问题 4 终点/h & 问题 4 变化/h \\
    \midrule
    末小时最低 $T_\infty$（$-0.408$ K） & 57.9143 & +0.7462 & 51.4786 & +0.6545 \\
    末小时最高 $T_\infty$（$+0.071$ K） & 57.0396 & -0.1285 & 50.7115 & -0.1127 \\
    末小时最低 $C_\infty$（$-1.1\times10^{-4}$ kg/kg） & 57.1659 & -0.0022 & 50.8217 & -0.0024 \\
    末小时最高 $C_\infty$（$+3.8\times10^{-4}$ kg/kg） & 57.1757 & +0.0076 & 50.8327 & +0.0085 \\
    \bottomrule
  \end{tabular}
  \par\smallskip\footnotesize 扰动幅度是附件 1 最后 1 h 内 $T_\infty$、$C_\infty$ 的实测最小值与最大值相对末值的偏移；参照算例取 $N_r=161$。温度方向的敏感度比浓度方向高约两个数量级。
\end{table}
